% =====================================================================
%  Analisi dei Capitolati  --  Git Push Team (Gruppo 9)
%  Documento esterno della Candidatura, diviso in capitoli:
%     analisi-dei-capitolati.tex   <- questo file (si compila questo)
%     capitoli/*.tex               <- un file per capitolo
% =====================================================================
\documentclass[documento]{gpt}

\titolo{Analisi dei Capitolati}
\uso{esterno}
\responsabile{spagnoli}

% Il capitolato in copertina si aggiunge solo dopo l'aggiudicazione
% (3 novembre 2026):
% \capitolato{C2 -- NEXUM Concierge (Eggon S.r.l.)}

% ---------------------------------------------------------------------
%  REGISTRO DELLE MODIFICHE  --  la riga più recente va IN CIMA.
%  Alla verifica finale il Verificatore aggiunge in cima la riga 1.0.0.
% ---------------------------------------------------------------------
\modifica{0.9.1}{2026-10-26}{guglielmo}{--}{Termini marcati allineati al Glossario}
\modifica{0.9.0}{2026-10-26}{guglielmo}{--}{Revisione complessiva: sintesi dei contenuti e approfondimento dei capitolati C3--C9}
\modifica{0.8.0}{2026-10-26}{biasini}{corro}{Integrazione degli esiti dell'incontro con Eggon del 26 ottobre}
\modifica{0.7.0}{2026-10-26}{guglielmo}{cescobolla}{Valutazione finale di C2, tabelle di confronto e conclusioni}
\modifica{0.6.0}{2026-10-25}{biasini}{corro}{Stesura della sintesi dei capitolati C3--C9}
\modifica{0.5.0}{2026-10-24}{guglielmo}{cescobolla}{Stesura dell'analisi del capitolato scelto, C2 -- NEXUM Concierge}
\modifica{0.4.0}{2026-10-23}{biasini}{corro}{Stesura dell'analisi di C1 -- AURA con gli esiti dell'incontro del 22 ottobre}
\modifica{0.3.0}{2026-10-22}{guglielmo}{cescobolla}{Stesura dell'analisi di C0 -- AIRA con gli esiti dell'incontro del 21 ottobre}
\modifica{0.2.0}{2026-10-21}{biasini}{corro}{Stesura della metodologia e dei criteri di valutazione}
\modifica{0.1.0}{2026-10-20}{guglielmo}{cescobolla}{Prima stesura: struttura del documento e introduzione}

% Tabella a due colonne per punti di forza e criticità:
%   \begin{fortedeboli} \fd{forza}{criticità} ... \end{fortedeboli}
\newenvironment{fortedeboli}{%
  \par\noindent\tabularx{\linewidth}{@{}>{\raggedright\arraybackslash}X
      >{\raggedright\arraybackslash}X@{}}
    \toprule
    \intestazione{Punti di forza} & \intestazione{Criticità} \\
    \midrule}
  {\bottomrule\endtabularx\par}
\newcommand{\fd}[2]{#1 & #2 \\[0.3em]}

% Scheda sintetica di un capitolato:  \begin{scheda} \voce{a}{b} \end{scheda}
\newenvironment{scheda}{%
  \par\noindent\tabularx{\linewidth}{@{}>{\raggedright\arraybackslash}p{0.24\linewidth}
      >{\raggedright\arraybackslash}X@{}}
    \toprule}
  {\bottomrule\endtabularx\par}
\newcommand{\voce}[2]{\textsf{\small\bfseries #1} & #2 \\}

% Indice fino alle sottosezioni
\setcounter{tocdepth}{2}

\begin{document}

% !TEX root = ../analisi-dei-capitolati.tex
% Capitolo 1 - Introduzione

\section{Introduzione}

\subsection{Scopo del documento}
Il documento riporta l'analisi dei \glo[capitolato]{capitolati} d'appalto del
progetto didattico di Ingegneria del Software 2026/2027 e motiva la scelta
del Git Push Team (Gruppo~9): il capitolato \textbf{C2 -- NEXUM
Concierge} di Eggon S.r.l. Sono analizzati in dettaglio anche gli altri
due capitolati della rosa finale, \textbf{C1 -- AURA} (Vimar S.p.A.) e
\textbf{C0 -- AIRA} (Sanmarco Informatica S.p.A.), mentre i restanti sono
descritti più brevemente.

I termini contrassegnati dalla lettera G a pedice, come \glo{RAG}, sono
spiegati nel Glossario.

\subsection{Riferimenti}
\begin{itemize}
  \item capitolati C0--C9 dell'anno accademico 2026/2027; per C2 il
        \emph{Business Requirements Document} BRD-AI-CONCIERGE-2026,
        versione~9 del 13 ottobre 2026;
  \item verbali interni VI~2026-10-12, VI~2026-10-13 e VI~2026-10-19;
  \item verbali esterni VE~2026-10-21 (Sanmarco Informatica),
        VE~2026-10-22 (Vimar) e VE~2026-10-26 (Eggon).
\end{itemize}

% !TEX root = ../analisi-dei-capitolati.tex
% Capitolo 2 - Criteri e metodologia

\section{Criteri e metodologia}\label{sec:metodologia}

\subsection{Criteri di valutazione}
Per confrontare proposte molto diverse il gruppo ha usato sei criteri:
\begin{itemize}
  \item \textbf{interesse}: coinvolgimento dei membri verso il dominio e le
        tecnologie;
  \item \textbf{valore formativo}: tecnologie e metodi spendibili dopo il
        corso, in un contesto vicino a un prodotto reale;
  \item \textbf{chiarezza del perimetro}: precisione di requisiti, criteri
        di accettazione e prodotti attesi;
  \item \textbf{sostenibilità}: equilibrio tra sfida tecnica e carico di
        lavoro di sei persone nel budget orario del progetto;
  \item \textbf{vincoli esterni}: hardware da gestire, servizi a pagamento,
        tecnologie imposte, integrazioni non controllate dal gruppo;
  \item \textbf{supporto del proponente}: incontri periodici, canali
        rapidi, materiale di prova.
\end{itemize}

\subsection{Processo decisionale}
\begin{enumerate}
  \item \textbf{Lettura individuale} (12--13 ottobre): ogni membro ha letto
        i capitolati e approfondito quelli di suo interesse.
  \item \textbf{Prima selezione} (13 ottobre): dopo la presentazione e la
        discussione, i più votati sono stati C0, C1 e C2; una seconda
        votazione tra i soli finalisti ha dato l'ordine di preferenza della
        Tabella~\ref{tab:votazione}.
  \item \textbf{Approfondimento a coppie}: C0 a \nomi{cescobolla,corro}, C1
        a \nomi{biasini,guglielmo}, C2 a \nomi{mazzurana,spagnoli}. Ogni
        coppia ha preparato le domande per il proponente, validate dal
        gruppo.
  \item \textbf{Incontri con i proponenti}: Sanmarco Informatica il 21
        ottobre, sessione congiunta di Vimar il 22 ottobre, Eggon il 26
        ottobre.
  \item \textbf{Decisione finale}: conferma di C2, con C1 e C0 come
        alternative nell'ordine.
\end{enumerate}

\begin{table}[H]
  \centering
  \caption{Votazione tra i capitolati finalisti (13 ottobre 2026)\label{tab:votazione}}
  \begin{tabular}{@{}llrr@{}}
    \toprule
    \intestazione{Capitolato} & \intestazione{Proponente} &
    \intestazione{Voti} & \intestazione{Percentuale} \\
    \midrule
    \textbf{C2} -- NEXUM Concierge & Eggon S.r.l.                & 3 & 50\% \\
    \textbf{C1} -- AURA            & Vimar S.p.A.                & 2 & 33\% \\
    \textbf{C0} -- AIRA            & Sanmarco Informatica S.p.A. & 1 & 17\% \\
    \bottomrule
  \end{tabular}
\end{table}

% !TEX root = ../analisi-dei-capitolati.tex
% Capitolo 3 - Capitolato scelto: C2 - NEXUM Concierge

\section{Capitolato scelto: C2 -- NEXUM Concierge}\label{sec:c2}

\begin{scheda}
  \voce{Proponente}{Eggon S.r.l., Padova}
  \voce{Referente}{Gianpaolo Ferrarin}
  \voce{Dominio}{assistente conversazionale per i servizi HR di una
    piattaforma esistente}
\end{scheda}

\subsection{Obiettivo}
Eggon sviluppa \textbf{NEXUM}, una piattaforma per la comunicazione
aziendale e i servizi HR composta da un \glo{backoffice} amministrativo e
da una \glo{PWA} per il personale. Il progetto chiede di introdurre nella
PWA un \textbf{Concierge}: un assistente conversazionale con cui il
personale trova informazioni, capisce le procedure e usa i servizi
aziendali in linguaggio naturale, mentre dal backoffice l'azienda gestisce
le fonti e le capacità dell'assistente.

Per restare affrontabile, l'\glo{MVP} è limitato a \textbf{ferie e
permessi}, ma l'architettura deve permettere di aggiungere in futuro
altre fonti e abilità senza riprogettare il sistema. Il capitolato
prosegue la collaborazione avviata con l'Università nel 2025/2026, ma il
nuovo progetto ne è indipendente.

\subsection{Caratteristiche funzionali}
Il cuore del capitolato è riconoscere quattro tipi di richiesta, anche
all'interno della stessa conversazione (Tabella~\ref{tab:c2-capacita}).

\begin{table}[H]
  \centering
  \caption{Le quattro capacità del Concierge\label{tab:c2-capacita}}
  \begin{tabularx}{\linewidth}{@{}>{\raggedright\arraybackslash}p{3.3cm}
      >{\raggedright\arraybackslash}X >{\raggedright\arraybackslash}X@{}}
    \toprule
    \intestazione{Capacità} & \intestazione{Esempio} &
    \intestazione{Fonte o strumento} \\
    \midrule
    Conoscenza aziendale & «Come funzionano le ferie in azienda?» &
      Regolamenti caricati dal backoffice \\
    Conoscenza del prodotto & «Dove posso segnare le ferie?» & Manuali di
      NEXUM \\
    Dati personali & «Quante ferie ho fatto quest'anno?» & Abilità di
      lettura dei dati \\
    Operazioni & «Registra le mie ferie dal 10 al 15 ottobre» & Abilità di
      registrazione della richiesta \\
    \bottomrule
  \end{tabularx}
\end{table}

Attorno a queste capacità il BRD richiede:
\begin{itemize}
  \item una \textbf{chat testuale} che mantiene il contesto e chiede
        chiarimenti quando la richiesta è ambigua (per esempio una data
        senza anno);
  \item un \textbf{backoffice} per caricare, sostituire e ritirare le fonti
        (inizialmente PDF), organizzate in \glo{bucket} con destinatari
        autorizzati; le risposte devono citare le fonti e il sistema deve
        astenersi quando le informazioni non bastano;
  \item \textbf{dati personali} letti solo tramite i servizi di NEXUM: il
        modello non deve inventare dati né fare conteggi al posto dei
        servizi;
  \item un \textbf{riepilogo completo e una conferma esplicita} prima di
        ogni invio, distinguendo tra richiesta registrata e approvata;
  \item \textbf{autorizzazioni, gestione degli errori e tracciabilità}:
        due profili (personale e amministrazione), separazione tra aziende,
        registrazione delle fonti usate e degli esiti delle operazioni.
\end{itemize}

\subsection{Requisiti e criteri di accettazione}
Il BRD elenca sedici requisiti di business, che saranno la base
dell'Analisi dei Requisiti. Non distingue requisiti obbligatori e
opzionali: il nucleo obbligatorio è l'intero percorso sul dominio scelto,
mentre grafi di conoscenza, distillazione avanzata, altri formati e altri
domini sono possibilità ulteriori. Per l'accettazione l'MVP deve
dimostrare tutte e quattro le capacità e un percorso unitario dalla
domanda informativa all'invio della richiesta, verificando inoltre
chiarimenti, citazioni e astensione, aggiornamento dopo il ritiro di una
fonte, autorizzazioni e gestione degli errori.

Il lavoro è diviso in tre fasi: un \glo{PoC} con le interfacce concordate
(alla RTB), poi il percorso completo e la verifica dell'MVP (alla PB).
L'MVP può usare servizi NEXUM di prova o simulatori conformi al contratto
di integrazione definito da Eggon, purché il percorso sia eseguibile da
capo a fondo.

\subsection{Considerazioni tecnologiche}
Il capitolato non impone tecnologie. Nell'incontro Eggon ha indicato come
riferimento non vincolante il proprio stack: Ruby on Rails per il backend,
Angular per backoffice e PWA, RubyLLM come interfaccia unica verso i
provider di modelli (per esempio Amazon Bedrock, Google Vertex,
OpenRouter), infrastruttura AWS. Chiede però un approccio
\textbf{agnostico rispetto al provider} e una scelta motivata da
esperimenti:
\begin{itemize}
  \item per la \textbf{conoscenza documentale}, confrontare il documento
        inserito nel contesto della conversazione con soluzioni basate su
        \glo{embedding}, ricerca vettoriale ed eventualmente grafi
        (\glo{RAG});
  \item per le \textbf{operazioni}, esporre le abilità tramite \glo{MCP},
        con un registro che ne definisce ingressi e uscite e un
        classificatore che riconosce l'intenzione dell'utente (posizione
        ferie, richiesta, conferma, annullamento, spostamento).
\end{itemize}
Il gruppo partirà da servizi simulati per poi passare ai modelli reali su
AWS, controllando il consumo di token.

\subsection{Incontro con il proponente}
Dall'incontro del 26 ottobre 2026 (VE~2026-10-26) sono emersi questi
chiarimenti:
\begin{itemize}
  \item \textbf{perimetro}: MVP solo su ferie e permessi, con il
        regolamento aziendale in PDF come primo documento e solo lingua
        italiana;
  \item \textbf{ricerca applicata}: il valore atteso sta nel confronto tra
        approcci e nella qualità delle decisioni, non solo nel PoC;
  \item \textbf{risorse}: accesso ad AWS e ai modelli con limiti di
        budget, i due repository dei gruppi precedenti come riferimento,
        esempi e materiali di prova costruiti insieme;
  \item \textbf{dati}: solo dati fittizi in sperimentazione;
  \item \textbf{gruppo}: sei persone invece delle sette del BRD non sono un
        problema;
  \item \textbf{collaborazione}: incontro settimanale di circa un'ora in un
        giorno fisso, Telegram per le comunicazioni rapide, Meet per gli
        incontri; Eggon chiede un confronto proattivo sulle scelte, anche
        nei periodi di esami;
  \item \textbf{strumenti di generazione del codice}: ammessi, purché il
        gruppo sappia spiegare il codice e ne curi la sicurezza (OWASP,
        Cyber Resilience Act); Eggon è disponibile a mostrare come lavora
        con Claude Code.
\end{itemize}

\subsection{Rischi principali}
\begin{itemize}
  \item \textbf{Ampliamento del perimetro}: mantenere ferie e permessi come
        unico dominio dell'MVP.
  \item \textbf{Risposte non supportate dalle fonti}: citazioni
        verificabili, astensione e un insieme di domande di prova per
        misurare la qualità.
  \item \textbf{Invii errati o duplicati}: riepilogo, conferma esplicita e
        verifica dell'esito prima di ripetere un'operazione.
  \item \textbf{Dipendenza dalle integrazioni}: concordare presto il
        contratto con Eggon e sviluppare su simulatori conformi.
  \item \textbf{Tecnologie nuove}: studio mirato nelle prime settimane e
        incontri di approfondimento con il proponente.
  \item \textbf{Continuità della comunicazione}, indicata da Eggon come
        rischio principale: incontro settimanale fisso anche nei periodi di
        esami.
  \item \textbf{Costi dei modelli}: servizi simulati all'inizio e
        monitoraggio del consumo di token entro il budget concordato.
\end{itemize}

\subsection{Valutazione del gruppo}

\begin{fortedeboli}
  \fd{\textbf{Perimetro circoscritto}: un solo dominio e vincoli espliciti
      su ciò che non è richiesto.}
     {\textbf{Scelte da sperimentare}: stack e strategia documentale vanno
      decisi e motivati con esperimenti.}
  \fd{\textbf{Criteri di accettazione verificabili}: ottima base per
      requisiti e test.}
     {\textbf{Qualità delle risposte difficile da misurare}: servono
      metriche e domande di prova costruite dal gruppo.}
  \fd{\textbf{Tecnologie attuali}: RAG, MCP e uso di strumenti da parte
      di un \glo{LLM}.}
     {\textbf{Tecnologie nuove per il gruppo}: serve un investimento
      iniziale di studio.}
  \fd{\textbf{Prodotto reale e proponente presente}: infrastruttura
      fornita, incontri settimanali, materiale dei gruppi precedenti.}
     {\textbf{Dipendenze}: contratto di integrazione da Eggon e budget
      per i modelli da rispettare.}
\end{fortedeboli}

Il gruppo ha scelto C2 per l'equilibrio tra interesse, valore formativo e
sostenibilità: un tema attuale in un prodotto reale, con un perimetro
limitato e criteri verificabili, senza hardware né tecnologie imposte.
L'incontro con Eggon ha confermato la scelta: il proponente fornisce le
risorse necessarie e cerca proprio una collaborazione basata sulla
sperimentazione e sul confronto delle scelte.

% !TEX root = ../analisi-dei-capitolati.tex
% Capitolo 4 - C1 - AURA

\section{C1 -- AURA}\label{sec:c1}

\begin{scheda}
  \voce{Proponente}{Vimar S.p.A.}
  \voce{Dominio}{domotica e servizi di intelligenza artificiale generativa}
\end{scheda}

\subsection{Obiettivo}
Vimar chiede un server \glo{MCP} che renda lo standard \textbf{KNX IoT 3rd
Party API}, con cui si comandano i suoi impianti domotici, utilizzabile dai
principali servizi commerciali di intelligenza artificiale (Claude,
ChatGPT, Copilot, Gemini). L'utente può così interrogare e comandare la
propria casa in linguaggio naturale con l'assistente che preferisce.

\subsection{Requisiti principali}
Sono obbligatori tre componenti: un'infrastruttura a container descritta
come codice con Docker Compose e compatibile con AWS; un server MCP con
autenticazione OAuth2, selezione dell'impianto, lettura e scrittura sui
dispositivi, più un client MCP per i test; il collegamento a un servizio
di intelligenza artificiale commerciale, che deve gestire 17 casi d'uso
base e 21 avanzati. I test devono coprire almeno il 75\% del codice e
l'80\% dei requisiti end-to-end. Sono opzionali il rilascio su AWS, il
client via VS Code, CLI o TUI e la copertura oltre il 90\%; storico,
scenari composti e trigger sono \emph{nice-to-have}. Vimar fornisce un kit
di impianto fisico per gruppo e consiglia Python con l'SDK FastMCP.

\subsection{Incontro con il proponente}
Alla sessione congiunta del 22 ottobre 2026 (VE~2026-10-22) Vimar ha
chiarito che il linguaggio è libero, che la soluzione è principalmente
back-end e basta che giri in locale con Docker, che per lo stato il
\emph{polling} è accettato, che il kit resta a casa di un membro e si
comanda via cloud, e che il PoC minimo copre due o tre casi base. Nei test
il servizio di intelligenza artificiale va simulato, perché un LLM reale
non è deterministico.

\subsection{Valutazione del gruppo}

\begin{fortedeboli}
  \fd{\textbf{Tema innovativo}: MCP collegato a un impianto reale.}
     {\textbf{Hardware}: il kit fisico condiviso è un punto singolo di
      guasto.}
  \fd{\textbf{Requisiti molto dettagliati} e proponente molto
      strutturato.}
     {\textbf{Dominio da studiare}: KNX IoT e documentazione per
      integratori.}
  \fd{\textbf{Materiale fornito}: kit, API ed eventuali risorse AWS.}
     {\textbf{Parte conversazionale delegata} a servizi esterni.}
\end{fortedeboli}

C1 è la prima alternativa del gruppo: solido e ben seguito, ma con la
dipendenza da un dispositivo fisico e un interesse leggermente inferiore
per un progetto in cui l'intelligenza artificiale è affidata a servizi
commerciali.

% !TEX root = ../analisi-dei-capitolati.tex
% Capitolo 5 - C0 - AIRA

\section{C0 -- AIRA}\label{sec:c0}

\begin{scheda}
  \voce{Proponente}{Sanmarco Informatica S.p.A.}
  \voce{Dominio}{automazione delle richieste B2B multicanale}
\end{scheda}

\subsection{Obiettivo}
Le richieste dei clienti B2B (ordini, ticket, appuntamenti, documenti
contabili) arrivano da canali eterogenei: e-mail con allegati, note
vocali, messaggi, foto di cataloghi. Il progetto chiede uno strumento che
monitori i canali, interpreti le richieste con trascrizione vocale,
\glo{OCR} e un LLM, ne estragga dati strutturati e li inoltri in
automatico ai gestionali aziendali tramite API REST o e-mail.

\subsection{Requisiti principali}
Sono obbligatori il controllo continuo dei canali, un filtro di sicurezza
in ingresso, l'elaborazione semantica tramite LLM, l'avvio di azioni
configurabili e un'interfaccia amministrativa web; sono opzionali una
dashboard di monitoraggio e le risposte automatiche al mittente. Le
tecnologie consigliate sono Angular, Java con Spring Boot, LangChain4j,
Whisper o Tesseract per voce e immagini e un database relazionale.

\subsection{Incontro con il proponente}
All'incontro del 21 ottobre 2026 (VE~2026-10-21) Sanmarco Informatica ha
chiarito che non ci saranno costi esterni (modelli \emph{open weight} in
locale, per esempio Qwen), che WhatsApp e i gestionali saranno simulati e
che i canali prioritari sono e-mail e audio. Ha inoltre chiesto misure di
sicurezza oltre la whitelist: antispam, offuscamento e cifratura dei dati,
isolamento dei tenant. Il supporto prevede un incontro ogni due settimane,
poi settimanale dopo il PoC.

\subsection{Valutazione del gruppo}

\begin{fortedeboli}
  \fd{\textbf{Problema concreto} e diffuso nelle aziende.}
     {\textbf{Ampiezza tecnologica}: voce, OCR, LLM, integrazioni e
      interfaccia web insieme.}
  \fd{\textbf{Nessun costo}: modelli locali e canali simulati.}
     {\textbf{Limiti dei modelli locali}, legati all'hardware dei membri.}
  \fd{\textbf{Priorità chiarite} e supporto regolare.}
     {\textbf{Sicurezza estesa} su un nucleo già ampio.}
\end{fortedeboli}

C0 è la seconda alternativa del gruppo: concreto e senza costi, ma il più
ampio dei tre finalisti sul piano tecnologico e il meno votato.

% !TEX root = ../analisi-dei-capitolati.tex
% Capitolo 6 - Altri capitolati (C3-C9)

\section{Altri capitolati}\label{sec:altri}

I capitolati C3--C9 sono stati letti e discussi da tutto il gruppo nella
fase di prima selezione, ma non sono entrati tra i finalisti e non sono
stati approfonditi con i proponenti. Per ciascuno si riportano obiettivo,
requisiti e tecnologie principali e la valutazione del gruppo.

% ---------------------------------------------------------------------
\subsection{C3 -- FossDot}
\textbf{Proponente}: Uranion EMEA Ltd.

\textbf{Obiettivo.} Gli smart speaker in commercio sono dispositivi chiusi,
legati al cloud del produttore e a un solo assistente vocale. Uranion
propone un ecosistema open source completo per un assistente vocale
domestico sul modello dell'Echo Dot, in cui il proprietario sceglie il
provider di intelligenza artificiale (in cloud o locale), collega servizi
esterni tramite \glo{MCP} e fa cooperare più dispositivi in casa.

\textbf{Requisiti e tecnologie.} L'MVP comprende la scelta e la validazione
dell'hardware di riferimento (scheda, array di microfoni, altoparlante) con
distinta base e prototipo, un firmware che cattura l'audio, riconosce la
parola di attivazione e trasmette la voce al server, un servizio di
\emph{bridging} con pipeline speech-to-text, LLM e text-to-speech, e
un'interfaccia web per registrare i dispositivi. La consegna completa
aggiunge più provider commutabili, integrazione MCP, \emph{skill} e
agenti, funzioni multi-dispositivo e gestione sicura dei segreti, con una
copertura dei test software dell'80\%. Lo stack è libero, con licenza
Apache~2.0 o MIT; Rust è preferito per il firmware, e sono suggeriti
Raspberry Pi, ReSpeaker, openWakeWord, faster-whisper, Piper e Ollama.

\textbf{Valutazione.} Il tema della privacy e del controllo sui propri dati
è stato apprezzato. Il progetto richiede però competenze di elettronica e
di programmazione embedded lontane dal percorso del gruppo, l'acquisto e la
gestione di hardware e una dimostrazione finale su dispositivo fisico: un
rischio ritenuto non sostenibile.

% ---------------------------------------------------------------------
\subsection{C4 -- Piattaforma per la gestione dei veicoli e dei servizi post-vendita}
\textbf{Proponente}: M31 S.r.l.

\textbf{Obiettivo.} Le informazioni su modelli, veicoli venduti,
manutenzioni e campagne di richiamo sono sparse tra sistemi diversi. M31
chiede una piattaforma distribuita che colleghi casa automobilistica,
concessionarie e proprietari, con accesso separato ai dati in base al ruolo
e all'organizzazione.

\textbf{Requisiti e tecnologie.} Sono obbligatori la gestione di modelli,
veicoli, concessionarie e proprietari, le richieste di manutenzione con
aggiornamento dello stato, le campagne di richiamo con individuazione dei
veicoli coinvolti, le notifiche, tre interfacce web distinte e \glo{API}
sicure, documentate e versionate; i dati dei veicoli possono essere
inseriti a mano o simulati. Tra gli opzionali figurano storico tecnico,
passaggio di proprietà, preventivi, manutenzione predittiva e ricerca
geografica. Sono richiesti anche un'analisi dello stato dell'arte e un
\emph{test book}. Le tecnologie indicate sono microservizi in Node.js con
Nest.js e TypeScript, Go per le componenti più esigenti, NATS o Kafka per
gli eventi, Kubernetes su Google Cloud Platform, PostgreSQL, MongoDB,
Redis, Angular, Prometheus e Grafana.

\textbf{Valutazione.} Il capitolato è chiaro e ben strutturato, ma la parte
applicativa è un gestionale tradizionale e la complessità sta soprattutto
nell'infrastruttura distribuita (microservizi, code di eventi,
orchestrazione). Il gruppo ha preferito un progetto in cui l'innovazione
riguardi il comportamento del prodotto.

% ---------------------------------------------------------------------
\subsection{C5 -- RareCare}
\textbf{Proponente}: Miriade S.r.l.

\textbf{Obiettivo.} Le famiglie di persone con malattie rare devono
coordinarsi da sole con centri clinici, pediatri, terapisti, scuola e
servizi sociali. RareCare Navigator, pensato inizialmente per la Sindrome
di Sanfilippo, vuole affiancare la figura del \emph{Care Navigator} con un
assistente basato sull'intelligenza artificiale, un sistema di scadenze e
un archivio condiviso dei documenti di cura.

\textbf{Requisiti e tecnologie.} Sono obbligatori un modello di ruoli e
permessi granulare, un motore di scadenze e notifiche (rinnovi, piani
assistenziali ed educativi) e un'interfaccia conversazionale \glo{RAG} che
risponde solo sulla base di fonti ufficiali verificate; sono opzionali un
assistente vocale, la mappatura dei servizi territoriali e l'analisi
predittiva dei bisogni. Il nucleo della piattaforma deve essere
indipendente dalla patologia, con pacchetti di configurazione per
estenderla ad altre malattie. Sono raccomandati un'architettura serverless
su AWS (API Gateway, Lambda, Bedrock), un'app multipiattaforma in Flutter o
React Native, cifratura dei dati, autenticazione a più fattori e
accessibilità WCAG~2.1~AA.

\textbf{Valutazione.} Il valore sociale del progetto è stato riconosciuto
da tutto il gruppo e la componente RAG è affine a quella di C2. Hanno
pesato però la delicatezza del dominio (dati sanitari di minori,
conformità GDPR, possibili implicazioni regolatorie), l'ampiezza delle
responsabilità richieste e uno sviluppo mobile multipiattaforma che il
gruppo non intendeva affrontare.

% ---------------------------------------------------------------------
\subsection{C6 -- Caching semantico e dashboard FinOps per LiteLLM}
\textbf{Proponente}: Var Group S.p.A.

\textbf{Obiettivo.} Var Group usa il progetto open source LiteLLM come
orchestratore davanti a più provider di modelli linguistici. In produzione
ne ha riscontrato due limiti: la cache riconosce solo richieste identiche,
quindi domande equivalenti formulate in modo diverso vengono pagate due
volte, e non esiste una visione dei costi per team, progetto o utente.

\textbf{Requisiti e tecnologie.} Il progetto chiede di estendere LiteLLM
con una cache semantica basata su \glo{embedding} e ricerca per
similarità, con soglia configurabile e invalidazione, e con un servizio
FinOps che attribuisca i costi, gestisca budget e avvisi e preveda la
spesa mensile, con isolamento dei dati tra team. Prima dello sviluppo sono
richiesti una sessione di design thinking, lo user story mapping e un PoC
di fattibilità; i test devono coprire almeno il 70\% del codice. Le
tecnologie consigliate sono Python con FastAPI, Docker Compose,
sentence-transformers o gli embedding di OpenAI, Redis, Qdrant o
pgvector, PostgreSQL e React con Chart.js.

\textbf{Valutazione.} Il tema è attuale e ben delimitato, ma molto
orientato all'infrastruttura: il prodotto non ha un utente finale diretto
e il suo valore si misura soprattutto in risparmi di costo. L'interesse
del gruppo è stato inferiore rispetto ai finalisti.

% ---------------------------------------------------------------------
\subsection{C7 -- SmartDashboard}
\textbf{Proponente}: Ergon Informatica S.r.l.

\textbf{Obiettivo.} I dati delle aziende di produzione e distribuzione
restano spesso chiusi nel gestionale e ogni nuovo report richiede
l'intervento del reparto IT. Ergon chiede una piattaforma di
\glo{Business Intelligence} self-service che recuperi i dati dal
gestionale, li consolidi in un \glo{data warehouse} e permetta agli utenti
di costruire da soli le proprie dashboard.

\textbf{Requisiti e tecnologie.} La web app deve offrire un costruttore di
dashboard a trascinamento con grafici, KPI, tabelle pivot, gauge e mappe,
filtri trasversali, esportazione in PDF ed Excel e gestione degli utenti.
Un livello di intelligenza artificiale deve generare visualizzazioni da
richieste in linguaggio naturale e, come estensioni progressive, produrre
insight automatici, suggerire widget e spiegare gli indicatori. Il
proponente consiglia di costruire prima la pipeline di Business
Intelligence e poi il livello AI, e può fornire i dati di un caso di studio
reale. Le tecnologie suggerite sono DuckDB per il data warehouse (con
supporto desiderabile ad altri database), PostgreSQL o MariaDB, Polars per
le trasformazioni, React, Angular o Vue e un LLM commerciale.

\textbf{Valutazione.} Il capitolato è chiaro e propone un percorso
incrementale ragionevole. La parte prevalente del lavoro riguarda però la
pipeline dei dati e il costruttore di dashboard, ambiti che hanno
suscitato meno interesse rispetto agli assistenti conversazionali dei
finalisti.

% ---------------------------------------------------------------------
\subsection{C8 -- Sonar}
\textbf{Proponente}: Halue S.r.l. Società Benefit.

\textbf{Obiettivo.} Con gli anni uno store Shopify accumula app, webhook,
campi personalizzati, script e codice di tema di cui nessuno ricorda più
lo scopo, e questa conoscenza si perde a ogni cambio di fornitore. Sonar
deve collegarsi a uno store, acquisirne configurazioni e codice e
ricostruire in automatico cosa è installato, cosa reagisce a quali
eventi, quali campi sono davvero usati e cosa è orfano.

\textbf{Requisiti e tecnologie.} Le scansioni sono lunghe e soggette ai
limiti di frequenza delle API di Shopify: devono essere durevoli,
riprendibili e osservabili, e salvare ogni volta uno snapshot immutabile.
I dati confluiscono in un grafo di entità e dipendenze, su cui un motore di
regole estensibile segnala codice morto e configurazioni orfane con una
motivazione per ciascuna segnalazione; la documentazione è esportabile in
Markdown o HTML. Sono richiesti cifratura dei token, isolamento tra
organizzazioni, copertura dei test del 70\% e un \emph{test book}. Lo
stack è in gran parte vincolante: Next.js su Vercel, Vercel Workflows e
Cron Jobs, Supabase con autenticazione Google, TypeScript e Shopify Admin
GraphQL API.

\textbf{Valutazione.} È tra i capitolati più curati per requisiti e criteri
di accettazione, ma è legato a un ecosistema specifico con molti vincoli
tecnologici, e il dominio dell'e-commerce enterprise non ha coinvolto il
gruppo quanto i finalisti.

% ---------------------------------------------------------------------
\subsection{C9 -- SpecCheck}
\textbf{Proponente}: Zucchetti S.p.A.

\textbf{Obiettivo.} Con lo sviluppo assistito dagli LLM il collo di
bottiglia si è spostato dalla scrittura del codice alla sua verifica:
compilatori e test dicono se il codice funziona, non se fa ciò che era
stato chiesto. Zucchetti chiede di estendere l'agente di sviluppo
OpenCode con strumenti che verifichino la tracciabilità tra casi d'uso,
requisiti e codice sorgente.

\textbf{Requisiti e tecnologie.} Il prodotto è un pacchetto di skill,
comandi e almeno un plugin per OpenCode, installabile con uno script, che
legga i diagrammi e le descrizioni dei casi d'uso (per esempio in Mermaid,
XMI o Markdown) e le specifiche dei requisiti, li confronti tra loro e con
il comportamento del codice e mostri le discrepanze in una dashboard,
anche testuale. Il confronto finale, che richiede precisione, va
realizzato come plugin o programma esterno e non affidato solo al modello.
Sono opzionali una dashboard grafica, l'individuazione del codice non
coperto da specifiche e la lettura diretta da PDF. Il proponente fornisce
accesso a modelli di grandi dimensioni, con un tetto di spesa, e a
modelli più piccoli senza limiti.

\textbf{Valutazione.} Il tema è originale e molto vicino agli argomenti del
corso, ma il capitolato si dichiara fortemente esplorativo: parte dei
risultati dipende da capacità degli LLM non ancora consolidate. Il gruppo
ha preferito un progetto con perimetro e criteri di accettazione più
definiti.

% !TEX root = ../analisi-dei-capitolati.tex
% Capitolo 7 - Conclusioni

\section{Conclusioni}\label{sec:conclusioni}

La Tabella~\ref{tab:confronto} confronta i tre finalisti secondo i criteri
della Sezione~\ref{sec:metodologia}. \emph{Alto} è sempre il valore
favorevole: per i vincoli esterni indica vincoli contenuti.

\begin{table}[H]
  \centering
  \caption{Confronto tra i capitolati finalisti\label{tab:confronto}}
  \begin{tabularx}{\linewidth}{@{}Xccc@{}}
    \toprule
    \intestazione{Criterio} & \intestazione{C2 -- Concierge} &
    \intestazione{C1 -- AURA} & \intestazione{C0 -- AIRA} \\
    \midrule
    Interesse                & Alto  & Alto  & Medio \\
    Valore formativo         & Alto  & Alto  & Alto  \\
    Chiarezza del perimetro  & Alto  & Alto  & Medio \\
    Sostenibilità            & Medio & Medio & Basso \\
    Vincoli esterni          & Medio & Basso & Alto  \\
    Supporto del proponente  & Alto  & Alto  & Alto  \\
    \bottomrule
  \end{tabularx}
\end{table}

Il Git Push Team si candida per il capitolato \textbf{C2 -- NEXUM
Concierge} di Eggon S.r.l., che offre il miglior equilibrio tra
innovazione, perimetro chiaro e sostenibilità. In caso di mancata
aggiudicazione l'ordine di preferenza è \textbf{C1 -- AURA} e poi
\textbf{C0 -- AIRA}.


\end{document}
